La strategia piu' semplice e' l'unione delle segnalazioni: l'artefatto e' marcato come vulnerabile se almeno uno strumento segnala la famiglia $f$. In forma compatta,
\begin{equation}
s_{\mathrm{OR}}(x)=
\mathbb{1}\left[\sum_{i=1}^{L} e_i^f(x) \geq 1\right].
\end{equation}
Nel nostro esempio $\sum_i e_i=1$, quindi $s_{\mathrm{OR}}(x^\star)=1$ e la decisione e' vulnerabile. Questo comportamento e' utile come riferimento ad alta recall: se una CWE e' vista da un solo strumento, l'OR la conserva. Il costo e' altrettanto chiaro: ogni falso positivo di qualsiasi analizzatore entra direttamente nel risultato finale. Per questo motivo e' preferibile chiamarla ``unione'' o ``1-out-of-N'' invece di ``upper bound'': non e' un limite superiore della qualita', ma un limite superiore della copertura delle segnalazioni.

\begin{lstlisting}[caption={Implementazione OR in analysis/fusion/baselines.py}]
from analysis.fusion.baselines import or_predictions

preds = or_predictions(
    validation_df,
    tools,
    families,
    labels,
    fire_index,
)
\end{lstlisting}
